% =========================================================
% Proof: Monotonicity of Infimum
% Source: volume-iii/analysis/bounding/notes/notes-bounding-monotonicty.tex
% =========================================================

\subsection*{Monotonicity of Infimum}
\label{prf:monotonicity-of-infimum}

\begin{remark*}[Return]
$\leftarrow$ \hyperref[thm:monotonicity]{Back to Theorem in Notes}
\end{remark*}

\begin{proof}
\Claim Let $A, B \subseteq \mathbb{R}$ be nonempty and bounded below.
If $A \subseteq B$, then $\inf B \le \inf A$.

\Given $A \subseteq B$; both sets bounded below.
\Goal $\inf B \le \inf A$.

\ByThm{Completeness Axiom} $\inf B$ exists (as every nonempty bounded-below set
has a greatest lower bound).
\ByDef infimum, $\inf B \le b$ for all $b \in B$.

Since $A \subseteq B$, every $a \in A$ satisfies $a \in B$, hence $\inf B \le a$.
\Therefore $\inf B$ is a lower bound of $A$.

Since $\inf A$ is the greatest lower bound of $A$ and $\inf B$ is a lower
bound of $A$, we have $\inf B \le \inf A$. \AsReq
\end{proof}

\begin{remark*}[Proof shape]
Mirror image of the supremum monotonicity proof: identify $\inf B$ as a lower
bound of $A$, then invoke the greatest-lower-bound property of $\inf A$.
\end{remark*}

\begin{remark*}[Dependencies]
Completeness Axiom (existence of $\inf B$); definition of infimum.
\end{remark*}
